\section{Análisis y Discusión de Resultados}

La práctica permitió comprobar que la identificación de componentes y la medición de magnitudes eléctricas son procesos fundamentales dentro del trabajo experimental con circuitos electrónicos. A partir de las mediciones realizadas, se pudo contrastar el valor nominal indicado por el código de colores con el valor real medido con el multímetro, lo que permitió analizar el grado de precisión y la tolerancia de cada resistor. Las tres tablas que se muestran a continuación resumen esta información de manera sistemática: la primera presenta la comparación directa entre el valor nominal y el valor medido; la segunda muestra cómo se interpreta el código de colores para obtener el valor correcto; y la tercera ejemplifica la relación entre los valores nominales comunes y sus bandas correspondientes.

\subsection{Tabla 1: identificación de resistencias por código de colores y medición}

\begin{table}[H]
\centering
\caption{Comparación entre el valor nominal de una resistencia, la lectura experimental y la tolerancia.}
\renewcommand{\arraystretch}{1.2}
\resizebox{\textwidth}{!}{%
\begin{tabular}{|c|c|c|c|c|c|c|c|c|c|c|c|}
\hline
No. & Banda 1 & Banda 2 & Banda 3 & Banda 4 & Valor nominal & Valor medido & Tolerancia & $R_{\min} < R < R_{\max}$ & Error \% & Potencia (W) & $I_{\max}$ (mA) \\
\hline
0 & Marrón & Negro & Rojo & Dorado & 10 $\times 10^{2}$ $\Omega$ & 1020 $\Omega$ & 5\% & $950\,\Omega < R < 1050\,\Omega$ & 2\% & 1 & 31,6 \\
1 & Verde & Azul & Verde & Dorado & 5.6 $\times 10^{5}$ $\Omega$ & 5,57 M$\Omega$ & 5\% & $5{,}32\,\mathrm{k}\Omega < R < 5{,}88\,\mathrm{k}\Omega$ & 1\% & 0,25 & 0,2 \\
2 & Naranja & Naranja & Naranja & Dorado & 3,3 $\times 10^{3}$ $\Omega$ & 32,6 k$\Omega$ & 5\% & $31\,\mathrm{k}\Omega < R < 35\,\mathrm{k}\Omega$ & 1\% & 0,5 & 3,9 \\
3 & Amarillo & Violeta & Marrón & Dorado & 470 $\Omega$ & 460 $\Omega$ & 5\% & $437{,}5\,\Omega < R < 493{,}5\,\Omega$ & 2\% & 0,5 & 52,6 \\
4 & Marrón & Negro & Naranja & Dorado & 10 $\times 10^{3}$ $\Omega$ & 9,78 k$\Omega$ & 5\% & $9{,}5\,\mathrm{k}\Omega < R < 10{,}5\,\mathrm{k}\Omega$ & 2\% & 1 & 10 \\
5 & Marrón & Rojo & Rojo & Dorado & 1,2 $\times 10^{2}$ $\Omega$ & 1,180 k$\Omega$ & 5\% & $114\,\Omega < R < 126\,\Omega$ & 2\% & 2 & 40,8 \\
6 & Blanco & Marrón & Marrón & Dorado & 91 $\times 10^{1}$ $\Omega$ & 902 $\Omega$ & 5\% & $864{,}5\,\Omega < R < 955{,}5\,\Omega$ & 1\% & 0,5 & 23,4 \\
7 & Amarillo & Violeta & Naranja & Dorado & 47 $\times 10^{3}$ $\Omega$ & 47,3 k$\Omega$ & 5\% & $44{,}65\,\mathrm{k}\Omega < R < 49{,}35\,\mathrm{k}\Omega$ & 0,6\% & 0,25 & 2,3 \\
8 & Rojo & Rojo & Rojo & Dorado & 2,2 $\times 10^{3}$ $\Omega$ & 2,18 k$\Omega$ & 5\% & $2{,}09\,\mathrm{k}\Omega < R < 2{,}31\,\mathrm{k}\Omega$ & 1\% & 0,25 & 10,7 \\
9 & Rojo & Violeta & Verde & Dorado & 2,7 $\times 10^{5}$ $\Omega$ & 2,71 M$\Omega$ & 5\% & $256{,}5\,\mathrm{k}\Omega < R < 283{,}5\,\mathrm{k}\Omega$ & 1\% & 0,25 & 0,3 \\
10 & Naranja & Naranja & Rojo & Dorado & 3,3 $\times 10^{2}$ $\Omega$ & 3,3 k$\Omega$ & 5\% & $313{,}5\,\Omega < R < 346{,}5\,\Omega$ & 0\% & 0,5 & 12,3 \\
\hline
\end{tabular}
}%
\end{table}

Esta tabla muestra la relación entre el código de colores de cada resistencia, su valor nominal teórico y la lectura real registrada con el multímetro. Permite verificar si el componente se encuentra dentro del rango admisible definido por la tolerancia, y facilita el cálculo del error porcentual como una medida de precisión experimental. En la práctica, la mayoría de las resistencias se encontraron dentro del rango esperado, lo que confirma que la identificación por colores y la medición con el instrumento fueron consistentes.

\subsection{Tabla 2: interpretación del código de colores}

\begin{table}[H]
\centering
\caption{Conversión del código de colores a valor nominal y tolerancia.}
\renewcommand{\arraystretch}{1.2}
\begin{tabular}{|c|c|c|c|c|c|c|}
\hline
No. & Color 1 & Color 2 & Color 3 & Color 4 & Valor inicial & Notación correcta \\
\hline
0 & Marrón & Negro & Rojo & Dorado & $10\times10^{2}\Omega$ & 1 k$\Omega$ \\
1 & Rojo & Rojo & Naranja & Plateado & $22\times10^{3}\Omega$ & 22 k$\Omega$ \\
2 & Amarillo & Violeta & Verde & Dorado & $47\times10^{5}\Omega$ & 4,7 M$\Omega$ \\
3 & Verde & Marrón & Verde & Rojo & $51\times10^{5}\Omega$ & 5,1 M$\Omega$ \\
4 & Amarillo & Violeta & Negro & Dorado & $47\Omega$ & 47 $\Omega$ \\
5 & Marrón & Negro & Negro & Rojo & 10 $\Omega$ & 10 $\Omega$ \\
6 & Marrón & Negro & Dorado & Dorado & $10\times10^{0}\Omega$ & 1 $\Omega$ \\
7 & Rojo & Marrón & Rojo & Plateado & $21\times10^{2}\Omega$ & 2,1 k$\Omega$ \\
8 & Marrón & Rojo & Marrón & Dorado & $12\times10^{1}\Omega$ & 120 $\Omega$ \\
9 & Azul & Rojo & Negro & Plateado & 62 $\Omega$ & 62 $\Omega$ \\
10 & Amarillo & Violeta & Amarillo & Dorado & $47\times10^{3}\Omega$ & 470 k$\Omega$ \\
\hline
\end{tabular}
\end{table}

Esta segunda tabla muestra la forma en que se interpreta la secuencia de bandas de un resistor para obtener su valor nominal y su tolerancia. Cada color representa un número o un multiplicador, y la cuarta banda define el margen de error permitido. El análisis de esta tabla es útil porque permite comprender que la lectura del código de colores no es arbitraria, sino un sistema estándar que facilita la identificación rápida del componente en laboratorio.

\begin{table}[H]
\centering
\caption{Ejemplos de resistencias nominales y sus bandas de color correspondientes.}
\renewcommand{\arraystretch}{0,1}
\resizebox{\textwidth}{!}{%
\begin{tabular}{|c|c|c|c|c|c|}
\hline
No. & Valor nominal & Color 1 & Color 2 & Color 3 & Color 4 \\
\hline
0 & $1\,\mathrm{k}\Omega$ & Marrón & Negro & Rojo & Dorado \\
1 & $0{,}1\,\Omega$ & Marrón & Negro & Plateado & Dorado \\
2 & $10\,\Omega$ & Marrón & Negro & Negro & Dorado \\
3 & $10\,\mathrm{k}\Omega$ & Marrón & Negro & Naranja & Dorado \\
4 & $220\,\Omega$ & Rojo & Rojo & Marrón & Dorado \\
5 & $5\,\mathrm{M}\Omega$ & Verde & Negro & Verde & Dorado \\
6 & $100\,\mathrm{k}\Omega$ & Marrón & Negro & Amarillo & Dorado \\
7 & $1{,}2\,\mathrm{M}\Omega$ & Marrón & Rojo & Verde & Dorado \\
8 & $15\,\mathrm{k}\Omega$ & Marrón & Verde & Naranja & Dorado \\
9 & $1\,\Omega$ & Marrón & Negro & Dorado & Dorado \\
10 & $100\,\Omega$ & Marrón & Negro & Marrón & Dorado \\
\hline
\end{tabular}
}%
\end{table}

\subsection{Tabla 3: relación entre valores nominales comunes y bandas de color}


La tercera tabla permite relacionar directamente un valor nominal con las bandas de color que lo representan, lo cual es especialmente útil para comprobar la lectura del código de colores sin necesidad de depender exclusivamente del dato impreso en el cuerpo del componente. En términos prácticos, esta tabla evidencia que los valores nominales no se asignan al azar, sino que siguen un sistema lógico de cifras y multiplicadores que permite reconocer rápidamente la resistencia y su rango de tolerancia.

En conjunto, estas tres tablas muestran que el código de colores de una resistencia es un lenguaje técnico que permite identificar de manera rápida y confiable la magnitud del componente, mientras que la medición experimental valida la precisión del valor real con respecto al valor nominal esperado. Este análisis es esencial para interpretar correctamente los resultados del laboratorio y para garantizar que los componentes sean usados dentro de sus condiciones de operación adecuadas.
